% ****************************************************************************************************
% htwsaar-i-mst-config.tex 
% ****************************************************************************************************  
\RequirePackage[utf8]{inputenc}				
 \DeclareUnicodeCharacter{00A0}{~}								
\RequirePackage[T1]{fontenc} 
								
% ****************************************************************************************************
% 1. Personal data and user ad-hoc commands
% ****************************************************************************************************
\newcommand{\myTitle}{Evaluierung lokaler KI-Modelle auf Android für BCI-Anwendungen}
\newcommand{\myDegree}{Bachelor of Science (B.\,Sc.)\xspace}
%\newcommand{\myDegree}{Master of Science (M.\,Sc.)\xspace}
\newcommand{\myDegreeType}{Bachelor\xspace}
%\newcommand{\myDegreeType}{Master\xspace}
\newcommand{\myDegreeCourse}{Praktische Informatik}
%\newcommand{\myDegreeCourse}{Informatik und Web-Engineering}
%\newcommand{\myDegreeCourse}{Kommunikationsinformatik}
%\newcommand{\myDegreeCourse}{Produktionsinformatik}
\newcommand{\myName}{Ahmed Chaabo\xspace}
\newcommand{\myUni}{Hochschule für Technik und Wirtschaft des Saarlandes\xspace}
\newcommand{\myCompany}{SoftwareCenter Musterhausen\xspace}
\newcommand{\myFirstProf}{Prof. Dr. rer. nat. Dr. rer. med. habil. Daniel J.Strauss\xspace}
\newcommand{\mySecondProf}{Prof. Dr.-ing. Klaus Berberich\xspace}
\newcommand{\myLocation}{Saarbrücken\xspace}
\newcommand{\myTime}{Tag.~Monat~Jahr\xspace}
\newcommand{\currentVersion}{Version 2.25, August 2024\xspace} % TODO: ggf. über git Versionsinformationen automatisch bereitstellen und verwenden

% ********************************************************************
% Setup, finetuning, and useful commands
% ********************************************************************
\newcounter{dummy} % necessary for correct hyperlinks (to index, bib, etc.)
% ****************************************************************************************************


% ****************************************************************************************************
% 2. Loading some handy packages
% ****************************************************************************************************
% ******************************************************************** 
% Packages with options that might require adjustments
% ******************************************************************** 
\PassOptionsToPackage{ngerman}{babel}   % change this to your language(s)
 \RequirePackage{babel}					
 \RequirePackage{csquotes}
	
\PassOptionsToPackage{language=auto,style=numeric-comp,backend=biber,bibencoding=ascii,maxbibnames=50}{biblatex} % backend ggf. auf neueres Werkzeug (z.B. biber) anpassen
 \RequirePackage{biblatex}
 %\bibliography{Bibliography}	% alter Befehl
 \addbibresource{Bibliography.bib}

\PassOptionsToPackage{fleqn}{amsmath}		% math environments and more by the AMS 
 \RequirePackage{amsmath}

% ******************************************************************** 
% Setting up the page and margins
% ******************************************************************** 
\usepackage{geometry}
 \geometry{a4paper,left=25mm,right=35mm,top=25mm,bottom=30mm}
% DIESE WERTE SIND NICHT ZU VERÄNDERN -- DO NOT CHANGE THESE VALUES

% ******************************************************************** 
% General useful packages
% ******************************************************************** 
%\usepackage[automark]{scrpage2}
\PassOptionsToPackage{dvipsnames}{xcolor}
	\RequirePackage{xcolor} % [dvipsnames]  
	\definecolor{ingwi}{cmyk}{.9,0,0,0}
\usepackage{textcomp} % fix warning with missing font shapes
\usepackage{scrhack} % fix warnings when using KOMA with listings package          
\usepackage{xspace} % to get the spacing after macros right  
\usepackage{mparhack} % get marginpar right
%\usepackage{fixltx2e} % fixes some LaTeX stuff <-- ist seit 2015 nicht mehr notwendig
\PassOptionsToPackage{printonlyused}{acronym}
	\usepackage{acronym} % nice macros for handling all acronyms in the thesis
%\renewcommand{\bflabel}[1]{{#1}\hfill} % fix the list of acronyms
\usepackage{booktabs}
\usepackage{multirow}
\usepackage{todonotes} %Settings for ToDoNotes
% Eigene Shortcuts fuer laengere Befehle
	\newcommand{\todox}[1]{\todo[inline, size=\small]{#1}}
	%Nummerierte Anmerkungen
	\newcounter{todocounter}
	\renewcommand{\todox}[2][]{\stepcounter{todocounter}\todo[inline, size=\small,caption={\thetodocounter: #2}, #1]{\renewcommand{\baselinestretch}{0.5}\selectfont\thetodocounter: #2\par}}
\usepackage{blindtext}
%\usepackage{footmisc}
% ****************************************************************************************************


% ****************************************************************************************************
% 3. Setup floats: tables, (sub)figures, and captions
% ****************************************************************************************************
\usepackage{tabularx} % better tables
	\setlength{\extrarowheight}{3pt} % increase table row height
%\newcommand{\myfloatalign}{\centering} % to be used with each float for alignment
\usepackage{caption}
\captionsetup{format=hang,font=small}
\usepackage{subfig}
\usepackage{wrapfig}
\usepackage{tikz}
% ****************************************************************************************************


% ****************************************************************************************************
% 6. Setup code listings
% ****************************************************************************************************
\usepackage{listings} 
%\lstset{emph={trueIndex,root},emphstyle=\color{BlueViolet}}%\underbar} % for special keywords
\lstset{language=[LaTeX]Tex,%C++,
    keywordstyle=\color{RoyalBlue},%\bfseries,
    basicstyle=\small\ttfamily,
    %identifierstyle=\color{NavyBlue},
    commentstyle=\color{Green}\ttfamily,
    stringstyle=\rmfamily,
    numbers=none,%left,%
    numberstyle=\scriptsize,%\tiny
    stepnumber=5,
    numbersep=8pt,
    showstringspaces=false,
    breaklines=true,
    frameround=ftff,
    frame=single,
		texcl=true,
    belowcaptionskip=.75\baselineskip
    %frame=L
} 
%Styles für verschiedene Sprachen festlegen, z.B. Java
\lstdefinestyle{Java}{
belowcaptionskip=1\baselineskip,
  breaklines=true,
  xleftmargin=\parindent,
  language=Java,
	texcl=true,
  showstringspaces=false,
  basicstyle=\footnotesize\ttfamily,
  keywordstyle=\bfseries\color{green!40!black},
  commentstyle=\itshape\color{purple!40!black},
  identifierstyle=\color{blue},
  stringstyle=\color{orange}}

% Farben fuer den Quellcode-Stil (an gaengiger IDE-Optik orientiert)
\definecolor{lstframe}{gray}{0.6}         % grauer Rahmen
\definecolor{lstbg}{gray}{0.985}          % sehr heller Hintergrund
\definecolor{lstkw}{RGB}{20,40,150}       % Schluesselwoerter: dunkelblau
\definecolor{lsttype}{RGB}{0,128,128}     % Typen/Klassen: teal
\definecolor{lstconst}{RGB}{130,50,150}   % Konstanten: violett
\definecolor{lstnum}{RGB}{160,90,0}       % Zahlen: orange-braun
\definecolor{lstcomment}{gray}{0.5}       % Kommentare: grau
\definecolor{lststring}{RGB}{0,120,80}    % Strings: gruen

% Eigener Stil fuer Quellcode-Listings: Padding, kein Umbruch, nicht ueber Seiten splitten
\lstdefinestyle{code}{
  language=Java,
  basicstyle=\footnotesize\ttfamily\color{black!88},
  keywordstyle=\bfseries\color{lstkw},
  keywordstyle=[2]\color{lsttype},        % Typen/Klassen
  keywordstyle=[3]\color{lstconst},        % Konstanten
  commentstyle=\itshape\color{lstcomment},
  stringstyle=\color{lststring},
  showstringspaces=false,
  texcl=false,               % Kommentare literal (kein LaTeX) -> % und -> sind unproblematisch
  breaklines=false,          % kein Zeilenumbruch
  % zusaetzliche Wortklassen (2 = Typen, 3 = Konstanten)
  morekeywords=[2]{String,List,ArrayList,Arrays,SampleEntity,SessionFeatures,%
    Segment,FrequencyBand,Check,StringBuilder,Context,LlmInference,%
    LlmInferenceOptions,LlmEngine,StreamCallback,CountDownLatch,Metrics,%
    Runnable,Thread,View,Debug,MemoryInfo,BatteryManager,File,FileWriter,%
    EegFeatureExtractor,FeatureFormatter,PromptBuilder},
  morekeywords=[3]{NUM_BANDS,CHANNELS,SAMPLE_RATE,SEGMENT_COUNT,%
    MAX_INVALID_RATIO,OK,TOO_SHORT,TOO_MANY_INVALID},
  % Zahlen einfaerben
  literate=%
    {0}{{\color{lstnum}0}}1 {1}{{\color{lstnum}1}}1 {2}{{\color{lstnum}2}}1
    {3}{{\color{lstnum}3}}1 {4}{{\color{lstnum}4}}1 {5}{{\color{lstnum}5}}1
    {6}{{\color{lstnum}6}}1 {7}{{\color{lstnum}7}}1 {8}{{\color{lstnum}8}}1
    {9}{{\color{lstnum}9}}1,
  frame=single,
  rulecolor=\color{lstframe}, % grauer Rahmen statt schwarz
  backgroundcolor=\color{lstbg},
  framesep=8pt,              % Innenabstand: Rahmen -> Code (Padding)
  xleftmargin=14pt,          % Abstand links
  xrightmargin=14pt,         % Abstand rechts
  aboveskip=1.2\baselineskip,
  belowskip=1.2\baselineskip,
  belowcaptionskip=1\baselineskip,
  captionpos=b,
}
% ****************************************************************************************************


% ****************************************************************************************************
% 6. PDFLaTeX, hyperreferences and citation backreferences
% ****************************************************************************************************
% ********************************************************************
% Using PDFLaTeX
% ********************************************************************
\usepackage{xurl}  % better handling of long URLs / hyphenation
\PassOptionsToPackage{pdftex,hyperfootnotes=false,pdfpagelabels}{hyperref}
	\usepackage{hyperref}  % backref linktocpage pagebackref
\pdfcompresslevel=9
\pdfadjustspacing=1 
\PassOptionsToPackage{pdftex}{graphicx}
	\usepackage{graphicx} 
    

% ********************************************************************
% Hyperreferences
% ********************************************************************
\hypersetup{%
    %draft,	% = no hyperlinking at all (useful in b/w printouts)
    pdfstartpage=1, pdfstartview=Fit,%
	colorlinks=true, linktocpage=true,
	%urlcolor=Black, linkcolor=Black, citecolor=Black, %pagecolor=Black,%
	%urlcolor=brown, linkcolor=RoyalBlue, citecolor=green, %pagecolor=RoyalBlue,%
    % uncomment the following line if you want to have black links (e.g., for printing)
    colorlinks=false, pdfborder={0 0 0},
    breaklinks=true, pdfpagemode=UseNone, pageanchor=true, pdfpagemode=UseOutlines,%
    plainpages=false, bookmarksnumbered, bookmarksopen=true, bookmarksopenlevel=1,%
    hypertexnames=true, pdfhighlight=/O,%nesting=true,%frenchlinks,%
    pdftitle={\myTitle},%
    pdfauthor={\textcopyright\ \myName, \myUni},%
    pdfsubject={},%
    pdfkeywords={},%
    pdfcreator={pdfLaTeX},%
    pdfproducer={LaTeX with hyperref}%
}   

% ********************************************************************
% Setup autoreferences
% ********************************************************************
% There are some issues regarding autorefnames
% http://www.ureader.de/msg/136221647.aspx
% http://www.tex.ac.uk/cgi-bin/texfaq2html?label=latexwords
% you have to redefine the makros for the 
% language you use, e.g., american, ngerman
% (as chosen when loading babel/AtBeginDocument)
% ********************************************************************
\makeatletter
\@ifpackageloaded{babel}%
    {%
       \addto\extrasamerican{%
					\renewcommand*{\figureautorefname}{Figure}%
					\renewcommand*{\tableautorefname}{Table}%
					\renewcommand*{\partautorefname}{Part}%
					\renewcommand*{\chapterautorefname}{Chapter}%
					\renewcommand*{\sectionautorefname}{Section}%
					\renewcommand*{\subsectionautorefname}{Section}%
					\renewcommand*{\subsubsectionautorefname}{Section}% 	
				}%
       \addto\extrasngerman{% 
					\renewcommand*{\chapterautorefname}{Kapitel}%
					\renewcommand*{\sectionautorefname}{Abschnitt}%
					\renewcommand*{\subsectionautorefname}{Abschnitt}%
					\renewcommand*{\subsubsectionautorefname}{Abschnitt}% 
					\renewcommand*{\paragraphautorefname}{Absatz}%
					\renewcommand*{\subparagraphautorefname}{Absatz}%
					\renewcommand*{\footnoteautorefname}{Fußnote}%
					\renewcommand*{\FancyVerbLineautorefname}{Zeile}%
					\renewcommand*{\theoremautorefname}{Theorem}%
					\renewcommand*{\appendixautorefname}{Anhang}%
					\renewcommand*{\equationautorefname}{Gleichung}%        
					\renewcommand*{\itemautorefname}{Punkt}%
				}%	
			% Fix to getting autorefs for subfigures right (thanks to Belinda Vogt for changing the definition)
			\providecommand{\subfigureautorefname}{\figureautorefname}%  			
    }{\relax}
\makeatother


% ****************************************************************************************************
% 6. Last calls before the bar closes
% ****************************************************************************************************
% ********************************************************************
% Development Stuff
% ********************************************************************
%\listfiles
\PassOptionsToPackage{l2tabu,orthodox,abort}{nag}
	\usepackage{nag}
%\PassOptionsToPackage{warning, all}{onlyamsmath}
%	\usepackage{onlyamsmath}


% ****************************************************************************************************
% 7. Further adjustments (experimental)
% ****************************************************************************************************
%\usepackage{tocbibind} %Allows us to add Bibliography to ToC
\usepackage{enumitem}
\setdescription{font=\normalfont\bfseries} %Changes the appearance of description items
\usepackage[activate={true,nocompatibility},final,tracking=true,kerning=true,spacing=true,factor=1100,stretch=10,shrink=10]{microtype}


% ********************************************************************
% Using different fonts
% ********************************************************************
%\usepackage{lmodern} % <-- no osf support :-(
%\usepackage[tt=false]{libertine} % [osf]
%\usepackage{luximono} 
%\usepackage[urw-garamond]{mathdesign} <-- no osf support :-(
\usepackage{mathpazo} 
%\usepackage{libertinus} % math? osf?

\setkomafont{disposition}{\bfseries}
% ****************************************************************************************************
\usepackage{float}
\usepackage{tikz}
\usepackage{pdflscape}   % Querformat-Seiten (Rohdaten-Anhang)
\usepackage{longtable}   % seitenübergreifende Tabellen (Rohdaten-Anhang)

% ---- pgfplots (Ergebnis-Diagramme in der Evaluation) ----
\usepackage{pgfplots}
\pgfplotsset{compat=1.18}
\usepgfplotslibrary{groupplots}

\definecolor{mGemma1}{HTML}{1F77B4}   % Gemma 3 1B
\definecolor{mQwen}{HTML}{FF7F0E}      % Qwen 2.5 1.5B
\definecolor{mGemma2}{HTML}{7F7F7F}    % Gemma 2 2B

\pgfplotsset{
  eegbar/.style={
    ybar, bar width=7pt,
    width=6.8cm, height=5cm,
    enlarge x limits=0.3,
    symbolic x coords={Report, Artefakt, Vergleich},
    xtick=data, ymin=0, ymajorgrids=true,
    tick label style={font=\small}, label style={font=\small},
    legend style={font=\small, at={(0.5,1.18)}, anchor=south,
                  legend columns=3, draw=none, /tikz/every even column/.append style={column sep=8pt}},
  },
  eb/.style={error bars/y dir=both, error bars/y explicit},
}

% Screenshot mit abgerundeten Ecken und dezentem Rahmen
\newcommand{\screenshot}[2][\linewidth]{%
  \begin{tikzpicture}
    \node[inner sep=0pt, rounded corners=8pt, clip] (im) {\includegraphics[width=#1]{#2}};
    \draw[rounded corners=8pt, black!30, line width=0.5pt] (im.south west) rectangle (im.north east);
  \end{tikzpicture}%
}

\usepackage{pgfplots}
\pgfplotsset{compat=1.18}

\definecolor{mGemma1}{HTML}{1F77B4}   % Gemma 3 1B
\definecolor{mQwen}{HTML}{FF7F0E}      % Qwen 2.5 1.5B
\definecolor{mGemma2}{HTML}{7F7F7F}    % Gemma 2 2B

\pgfplotsset{
  eegbar/.style={
    ybar, bar width=7pt,
    width=6.8cm, height=5cm,
    enlarge x limits=0.3,
    symbolic x coords={Report, Artefakt, Vergleich},
    xtick=data, ymin=0, ymajorgrids=true,
    tick label style={font=\small}, label style={font=\small},
    legend style={font=\small, at={(0.5,1.18)}, anchor=south,
                  legend columns=3, draw=none, /tikz/every even column/.append style={column sep=8pt}},
  },
  eb/.style={error bars/y dir=both, error bars/y explicit},
}
\usepgfplotslibrary{groupplots}